\chapter{Cenários e Atores}

\begin{objetivos}
\begin{itemize}
\item Adicionar novos atores e escolher cenários;
\item Desenhar o seu próprio ator;
\item Organizar atores nas camadas do palco.
\end{itemize}
\end{objetivos}

\section{Adicionando atores}

No canto inferior direito da tela, o botão azul redondo abre o menu de \textbf{novo ator}. Você pode:

\begin{passos}
\item \textbf{Escolher} um ator da biblioteca (animais, pessoas, objetos, letras\ldots);
\item \textbf{Pintar} o seu próprio ator no editor de desenho;
\item \textbf{Surpresa:} o Scratch escolhe um ator para você;
\item \textbf{Enviar} uma imagem do computador;
\item \textbf{Tirar uma foto} com a câmera.
\end{passos}

\begin{dica}
Para apagar um ator, clique no pequeno \textbf{X} no canto da miniatura dele. Para dar um nome fácil, mude o campo \textbf{Ator} acima da lista.
\end{dica}

\section{Cada ator, o seu código}

\begin{cuidado}
O código fica no \textbf{ator selecionado}. Se você montar blocos com o gatinho selecionado, só o gatinho obedece! Para dar ordens a outro ator, clique nele na lista antes de montar os blocos.
\end{cuidado}

\section{Cenários}

O \textbf{cenário} é o fundo do palco. Ele é o ator especial chamado \textbf{Palco}, no canto inferior direito, que não anda nem gira, mas pode ter fantasias (cenários diferentes) e código próprio.

\begin{center}
\renewcommand{\arraystretch}{1.8}
\begin{tabular}{@{}ll@{}}
\bloco[scLooks]{mude o cenário para \menu{praia}} & troca o fundo \\
\bloco[scLooks]{próximo cenário} & passa para o fundo seguinte \\
\end{tabular}
\end{center}

\begin{vocabulario}[Vocabulário: cenário]
\textbf{Cenário} é a imagem de fundo do palco, onde a história acontece. Um projeto pode ter vários cenários e trocar entre eles.
\end{vocabulario}

\section{Desenhando o seu ator}

O \textbf{editor de desenho} tem ferramentas como pincel, linha, círculo, retângulo, texto e balde de tinta. Há dois modos:

\begin{description}[leftmargin=!,labelwidth=3.4cm,style=nextline]
\item[\textbf{Vetor}] O desenho é feito de formas que podem ser esticadas sem perder a qualidade. É o melhor para quem está começando.
\item[\textbf{Bitmap}] O desenho é feito de pontinhos (pixels), como uma pintura. Dá efeitos mais livres.
\end{description}

\begin{dica}
Clique no \textbf{ponto central} do desenho: é ele que define em torno de qual ponto o ator gira e cresce. Coloque-o no meio do ator!
\end{dica}

\section{Camadas}

Se dois atores ocupam o mesmo lugar, um aparece \textbf{na frente} do outro. Isso se chama \textbf{camada}. Para mudar, use:

\begin{center}
\renewcommand{\arraystretch}{1.8}
\begin{tabular}{@{}ll@{}}
\bloco[scLooks]{vá para a camada da frente} & o ator vai para o topo \\
\bloco[scLooks]{vá para trás \entrada{1} camadas} & o ator recua uma camada \\
\end{tabular}
\end{center}

\begin{missao}[Missão 5: Conversa entre amigos]
Adicione um segundo ator e faça os dois conversarem. O gatinho fala primeiro, o amigo responde.
\end{missao}

Para isso precisamos de um jeito de combinar \textbf{a vez de cada um}. Vamos usar o tempo: o amigo espera 2 segundos antes de responder.

\begin{codigo}[Código do gatinho]
\begin{pilha}
\bk[scEvents]{quando \bandeira\ for clicado}
\bk[scLooks]{diga \entrada{Oi! Tudo bem?} por \entrada{2} segundos}
\end{pilha}
\end{codigo}

\begin{codigo}[Código do amigo]
\begin{pilha}
\bk[scEvents]{quando \bandeira\ for clicado}
\bk[scControl]{espere \entrada{2} segundos}
\bk[scLooks]{diga \entrada{Tudo ótimo!} por \entrada{2} segundos}
\end{pilha}
\end{codigo}

\begin{curiosidade}
Existe um jeito mais seguro de combinar a vez de cada ator: as \textbf{mensagens}. Vamos aprender isso no capítulo sobre Eventos.
\end{curiosidade}

\section{Hora de praticar}

\begin{exercicio}[Novo cenário]
Escolha um cenário da biblioteca e faça o projeto \textbf{trocar de cenário} 1 segundo depois de a bandeira ser clicada.
\end{exercicio}

\begin{solucao}
Esse código vai no \textbf{Palco}, não no ator.
\begin{pilha}
\bk[scEvents]{quando \bandeira\ for clicado}
\bk[scControl]{espere \entrada{1} segundos}
\bk[scLooks]{próximo cenário}
\end{pilha}
\end{solucao}

\begin{exercicio}[Dois amigos]
Adicione um ator da biblioteca. Faça o gatinho andar 100 passos para a direita e o outro ator dizer ``Bem-vindo!'' depois de 1 segundo.
\end{exercicio}

\begin{solucao}
Código do gatinho:
\begin{pilha}
\bk[scEvents]{quando \bandeira\ for clicado}
\bk[scMotion]{mova \entrada{100} passos}
\end{pilha}
Código do outro ator:
\begin{pilha}
\bk[scEvents]{quando \bandeira\ for clicado}
\bk[scControl]{espere \entrada{1} segundos}
\bk[scLooks]{diga \entrada{Bem-vindo!} por \entrada{2} segundos}
\end{pilha}
\end{solucao}

\begin{exercicio}[Quem está na frente?]
Coloque o gatinho em cima de outro ator. Faça o gatinho \textbf{ficar atrás} dele quando a bandeira for clicada.
\end{exercicio}

\begin{solucao}
\begin{pilha}
\bk[scEvents]{quando \bandeira\ for clicado}
\bk[scLooks]{vá para trás \entrada{1} camadas}
\end{pilha}
\end{solucao}

\begin{desafio}[Minha pequena história]
Crie uma cena com \textbf{dois atores} e \textbf{dois cenários}. No primeiro cenário os atores conversam; depois o cenário muda e eles falam de novo. Desenhe um dos atores você mesmo!
\end{desafio}

\begin{solucao}
Sugestão de código do Palco:
\begin{pilha}
\bk[scEvents]{quando \bandeira\ for clicado}
\bk[scControl]{espere \entrada{5} segundos}
\bk[scLooks]{próximo cenário}
\end{pilha}
O ator 1 fala logo; o ator 2 espera 2 segundos para responder; e depois que o cenário mudar, use \bloco[scControl]{espere \entrada{6} segundos} e um novo \bloco[scLooks]{diga} em cada ator.
\end{solucao}
